You are given a sky map ``Map-OT01'' along with the question paper. This map
shows only stars but no diffuse objects.

% FIGURE NEEDED: sky map "Map-OT01" (stars only, no diffuse objects)

Four known stars (S1, S2, S3 and S4), which are present in the map above,
are given in the table below with their Common names, Bayer designations and
Equatorial coordinates.

\begin{center}
\begin{tabular}{cllll}
Sr. No. & Common Name & Bayer Name & RA & Dec \\
\hline
S1 & Alpheratz & $\alpha$ Andromedae & $00^h08^m24^s$ & $29^\circ 5'16''$ \\
S2 & Markab & $\alpha$ Pegasi & $23^h04^m46^s$ & $15^\circ 12'17''$ \\
S3 & Scheat & $\beta$ Pegasi & $23^h03^m47^s$ & $28^\circ 4'58''$ \\
S4 & Algenib & $\gamma$ Pegasi & $00^h13^m14^s$ & $15^\circ 10'59''$ \\
\end{tabular}
\end{center}

Complete the tasks (OT01.1) and (OT01.2) while you are planning the
observations.

\begin{description}
\item[(OT01.1)] Your first task is to mark these 4 stars (with a circle
  around each star) and label them as S1, S2, S3 and S4 on the sky map
  ``Map-OT01'' provided to you. \hfill[6]
\item[(OT01.2)] An astronomer has discovered a new diffuse object `Z' at the
  following coordinates ---
  RA: $21^h36^m10.6^s$, Dec: $-26^\circ 10'24.4''$.

  Mark the position of this diffuse object on the same sky map ``Map-OT01''
  with a $\oplus$ sign and label it as `Z'. Assume that a linear rectangular
  grid for equatorial coordinates is valid in the relevant region of the
  map. \hfill[7]
\end{description}

The following task is to be performed once you reach the telescope station.

On the screen diagonally opposite to your station, initially a welcome
message, followed by a sample sky (sky unrelated to the question) will be
displayed along with a countdown timer. You may use this time to orient the
telescope towards the screen and familiarize yourself with other equipment
provided at the station. At the end of this time a part of the sky given in
the map ``Map-OT01'' will be projected on the screen for the next 6 minutes.
Note that the scale of the projection shown on the screen is different from
the actual scale seen in the sky.

\begin{description}
\item[(OT01.3)] Find the new object `Z' with the telescope using any
  appropriate eyepiece. Then centre the object properly in the field of view
  of the eyepiece with the crosshair, and show it to the examiner at your
  station.

  At the end of 6 minutes, the projection will be blurred for 20 seconds. At
  this point you must step away from the telescope. The projection will be
  restored for the examiner to check the view through the telescope. This
  marks the end of the first task. \hfill[12]
\end{description}
